\documentclass[10pt,letterpaper,final]{report}
\usepackage[margin=.75in]{geometry}
\usepackage[utf8]{inputenc}
\usepackage{amsmath}
\usepackage{amsfonts}
\usepackage{amssymb}
\usepackage{amsthm}
\renewcommand\qedsymbol{$\blacksquare$}
\usepackage{enumerate}
\usepackage{enumitem}
\usepackage{hyperref}
\usepackage{pdfpages}
\usepackage{graphics}
\usepackage{graphicx}
\usepackage{tikz}
\usepackage{tikz-qtree}
\usepackage{forest}
\usepackage{float}
\usetikzlibrary{automata,arrows}
\usetikzlibrary{shapes.multipart, positioning}
\usepackage{mdframed}
\usepackage[
  backend=biber,
  style=numeric-comp,
  sorting=none
]{biblatex}
\addbibresource{references.bib}

\usepackage{minted}
\usemintedstyle{algol_nu}
\usepackage{xltabular}
\usepackage{pdflscape}
\usepackage{tcolorbox}
\usepackage{enumitem}
\usepackage{array}
\usepackage{todonotes}
\setuptodonotes{inline}  % render \todo notes as full-width blocks in the text flow, not margin boxes

% Based on template from COSC-417 at Towson University, originally authored by Marius Zimand

% To input a script, adapt the following:
% \inputminted[breaklines, breakanywhere, fontsize=\footnotesize, frame=single, fontsize=\footnotesize]{python}{script.py}

\begin{document}
\begin{center}
  \fbox{
    \vbox{
      \begin{flushleft}
        Jonathan Llovet \\  % authors' names
        EN.625.252.81 \\  %class
        2026-09-13 \\  % date
        Module 2 - Homework 2 \\ % ID
      \end{flushleft}
      \center{\Large{\textbf{Linear Algebra}}}
      %\end{mdframed}
    } % end vbox
  } % end fbox
  \vline
\end{center}

\medskip

\noindent\textbf{Directions}

Complete the following 4 problems showing all your work. You may use a calculator to check your work, but should write out (or TeX up) all the steps of your solution. Unless otherwise specified, you may skip some steps of row reduction with a calculator, but state that you did so. Please upload your work as a single .pdf file in the course.

In the problems below, let
\begin{align*}
  A & =
  \left[\begin{array}{rrr}
      1 & 1 & 0 \\
      1 & 2 & 1 \\
      0 & 1 & 1
    \end{array}\right]
\end{align*}

\medskip
\noindent\textbf{Problem 1:}
\medskip

Calculate rref($A$) by hand showing all the steps of your row reduction.

\medskip
\noindent\textbf{Answer:}
\medskip

\pagebreak

\noindent\textbf{Problem 2:}
\medskip

Using your answer to Problem 1, determine the general solution of $Ax = 0$.

\medskip
\noindent\textbf{Answer:}
\medskip

\pagebreak

\noindent\textbf{Problem 3:}
\medskip

Solve the matrix equation
$Ax =
  \left[\begin{array}{r}
      1 \\
      0 \\
      1
    \end{array}\right]$, or show that no solution exists.

\medskip
\noindent\textbf{Answer:}
\medskip

\pagebreak

\noindent\textbf{Problem 4:}
\medskip

Use 3D graphing software (e.g. Geogebra or Desmos) to give a geometric description of your answers to problems 2 and 3. (Note: even if a systems is inconsistent, you can still graph to show no common intersection occurring.)

\medskip
\noindent\textbf{Answer:}
\medskip

\pagebreak

\section*{Source Code}
The full source code, including LaTeX, tooling, other artifacts, and git history is available on my Github repository for the class here:
\begin{center}
  \url{https://github.com/jllovet/jhu-en-625-252-fa26-linear-algebra}
\end{center}

\printbibliography

\end{document}
